\documentclass[12pt]{report}

%---------------------------------------------------------------
% PACKAGES
%---------------------------------------------------------------
\usepackage[utf8]{inputenc}
\usepackage{setspace}
\usepackage{apacite}
\usepackage{graphicx}
\usepackage{amsmath}
\usepackage{amssymb}
\usepackage{lipsum}
\usepackage[margin=1in]{geometry}
%\usepackage[utf8]{inputenc}
\linespread{1.5}
\usepackage{ragged2e}
%\setcounter{secnumdepth}{0}
\usepackage[none]{hyphenat}
\setlength{\emergencystretch}{3em}
\usepackage{titlesec}
\usepackage{booktabs}

% --- Custom Chapter Format ---
\titleformat{\chapter}[hang]  {\normalfont\Huge\bfseries}   {Chapter \thechapter:} {1em} {}                
\titlespacing*{\chapter}  {0pt}     {0pt}   {2ex} 


%\onehalfspacing
%\setcounter{secnumdepth}{3}
\usepackage{natbib}
\usepackage{hyperref}
%---------------------------------------------------------------
\begin{document}
\Centering
\justifying
%\sloppy

%---------------------------------------------------------------
\title{Sentiment-driven volatility forecasting using GARCH and Random Forests}
\author{Elly Okinyo}
\date{2025}
\maketitle
\clearpage
\pagenumbering{roman}    % i, ii, iii, ...
\setcounter{page}{1}     % start from i
\tableofcontents
\newpage
\clearpage
\pagenumbering{arabic}   % 1, 2, 3, ...
\setcounter{page}{1}     % Chapter 1 starts on page 1
%===============================================================
\chapter{Introduction}
%===============================================================

\section{Background}

Bitcoin has grown from being a mere technological experiment when it was first introduced in 2008/09 to a becoming a leading globally traded financial asset with significant economic relevance. Since its introduction, Bitcoin has has had phases of retail speculation, high dealers' testing, and just recently, there has been mainstream adoption of its trade even at the national levels as is the case of El Salvador where it was adopted by the state as an official currency\citep{ArgenteVanPatten2024Bitcoin}. From 2017 through 2025, Bitcoin has been the market leader among cryptocurrencies and is constantly highly ranked in this class of assets as well as other top globally traded digital assets, with evolving regulatory developments as well as increasing liquidity which has resulted in enhanced adoption of Bitcoin as an asset.

Bitcoin exhibits high volatility which makes it different from typical investment assets, with the variation in daily returns frequently exceeding those observed in equities, commodities, and foreign exchange markets, with huge price swings often influenced by market sentiment, investor attention, regulatory announcements, and broader macro-financial developments. Given that this has become such an important asset, it raises the need to understand how to accurately forecast its volatility since Bitcoin holds highly significant value for both investors and regulators alike.

\section{Importance of volatility forecasting}

Volatility forecasting is an integral part of financial economics. The accurate predictions of conditional variance support in portfolio optimisation, hedging strategies, and risk management frameworks such as Value-at-Risk (VaR). For regulators and investors alike, volatility assessments enhance a deeper understanding of speculative bubbles, systemic risks, and market stability.

Conventional econometric models, especially the Autoregressive Conditional Heteroskedasticity (ARCH) model introduced by \cite{engle1982arch} and its generalised extension, GARCH, by \cite{bollerslev1986garch}, remain the foundation of volatility modelling in financial economics since there models capture core empirical regularities such as volatility clustering and fat-tailed returns. These models, however, run into obstacles brought about by the fact that cryptocurrencies often exhibit stronger persistence, asymmetry patterns distinct from equities, and non-linear, sentiment-driven fluctuations compared to the highly seasonal data points like the national accounting statistics.

\section{Research problem and research gap}

While ARCH and GARCH models have been widely applied to cryptocurrencies for example, by \cite{dyhrberg2016bitcoin}; \cite{chu2017garchcrypto}; \cite{katsiampa2017volatility}, several limitations remain. First of all, many studies rely on short or outdated datasets ending before key structural events, such as the approval of U.S. spot Bitcoin ETFs in 2024(\cite{Gensler2024SpotBitcoin}) which attracted more informed market participants to the Bitcoin spot market, possibly leading to a change in the composition of its investors, leading to a record net-inflows upwards of one and a half billion US dollars\citep{Babalos2025BitcoinETFs}. Secondly, few studies directly compare GARCH-family models with machine learning methods such as Random Forests (e.g., \cite{basher2022bitcoinRF}), especially in the context of sentiment-driven volatility. This clearly outlines a gap since investor behaviour has been demonstrated to influence greatly how Bitcoin performs as an asset\citep{katsiampa2017volatility}. This necessitates the need for a more accurate variable representing sentiment, hence the role of sentiment, proxied by Google Trends and other search-based indicators, remains understudied over long time spans. 

Although hybrid models combining GARCH and ML have shown strong performance (\cite{zahidetal2022hybridgarch}; \cite{han2025combinedmethod}), they largely remain unexplored in sentiment-augmented settings, hence as much as we understand clearly that these hybrid models improve quality, more assessment need to be explored on how these models would perform with the integration of sentiment.

From the foregoing,therefore, arises a need for a systematic comparison of GARCH-family models and Random Forests, with and without sentiment variables, over an extended period of time.

\section{Research objectives}

The study has four main objectives. First, I compare the forecasting accuracy of GARCH-family models, particularly GARCH(1,1) and IGARCH(1,1) and Random Forests using Bitcoin daily data. Second, the study evaluates the effect of incorporating sentiment, proxied by Google Trends, on volatility forecast accuracy. Next, I assess whether sentiment-augmented Random Forest models perform better than the traditional econometric approaches used to study volatility. Finally, the study aims to provide an updated evidence on Bitcoin volatility dynamics using a comprehensive dataset encompassing major regime shifts.


\section{The relevance of sentiment}

Cryptocurrency markets, and in particular, crypto price performances, are heavily influenced by behavioural factors, media narratives, online discussions, and investor attention. Prior studies such as \cite{mai2018silentmajority} and \cite{sapkota2022newsbitcoin} have shown that search intensity and social media conversations and trends can lead up to price volatility. Google trends which show the popularity of search on Google by analysing the sample of real, anonymised searches, shows interest over time, by region, and in relation to other topics. Therefore, trends in search for ``Bitcoin" is used to serve as a proxy for retail interest and attention shocks, which are especially relevant in cryptocurrency markets dominated by non-institutional traders. Integrating sentiment into volatility models, therefore, has the potential to capture otherwise unobservable market forces.

\section{Structure of the thesis}

The thesis proceeds as follows. chapter one introduces the background, motivation, research questions, and contribution. Chapter two then reviews theoretical foundations and empirical findings across econometric, behavioural, and machine-learning approaches. Chapter three describes the data, preprocessing steps, and modelling methodology (GARCH, Random Forests, and hybrid approaches). In chapter four, I present the  empirical results including descriptive statistics, model comparison, and forecasting performance. Chapter five provides an interpretation of results and how this compares to the previous literature, implications for stakeholders, and limitations. And finally, chapter six concludes with a summary and recommendations for future work.


\newpage

%===============================================================
\chapter{Literature review}
%===============================================================

\section{Introduction}
Over the past four decades, volatility forecasting has been a central part of financial econometrics, owing to the need to understand, quantify, and predict time-varying uncertainty in asset returns. With the emergence of Bitcoin and other digital assets, this research area has experienced renewed attention due to the unique behavioural, structural, and technological characteristics of cryptocurrency markets. Bitcoin, in particular, depicts extreme price swings, high-frequency sharp reversals, heavy-tailed return distributions, and heightened sensitivity to regulatory, technological, and behavioural influences, features which challenge traditional econometric models and motivate comparative evaluations involving machine learning and sentiment-enhanced approaches.

In this chapter, I provide a comprehensive, structured, and comprehensive review of the academic literature covering the foundational ARCH and GARCH theoretical canon, early cryptocurrency applications, asymmetry and regime dynamics, realised volatility models, behavioural and sentiment-driven forecasting, and the increasing use of machine learning and other hybrid models in the forecasting of volatility

Section 2.2 revisits the foundational ARCH and GARCH literature, outlining their theoretical principles and empirical implications. Section 2.3 discusses the first-generation cryptocurrency volatility studies that applied classical GARCH models to Bitcoin and other digital assets. Section 2.4 reviews contributions on asymmetry, smooth-transition dynamics, and realised volatility. Section 2.5 reviews behavioural, sentiment, and attention-based research. Section 2.6 examines machine learning and hybrid GARCH–ML models. Section 2.7 concludes by identifying the gaps addressed by this study.

\section{Foundational ARCH/GARCH canon (1982–1994)}
Financial return series often exhibit features inconsistent with classical homoskedastic linear models, and this necessitates the developments in the study of volatility. This inconsistency is often due to volatility clustering, heavy-tailed distributions, and serial dependence in higher-order moments. The foundational literature like the introduction of ARCH established the theoretical and empirical pillars that continue to support the contemporary financial econometric studies.

\subsection{ARCH models}

The seminal contribution of \cite{engle1982arch} introduced the Autoregressive Conditional Heteroskedasticity (ARCH) model which marked a major shift from constant-variance assumptions that had been prominent in prior literature. In Engle’s specification, the conditional variance is modelled as a linear function of lagged squared residuals and this structure captures volatility clustering in which periods of high volatility followed by high volatility and produces fat-tailed unconditional return distributions even under normally distributed innovations. Engle’s application to U.K. inflation demonstrated that macroeconomic uncertainty portrays strong temporal dependence and can be statistically modelled rather than treated as exogenous noise.

A number of key theoretical insights from the ARCH model came to the forefront. First of all, the comparison of the conditional versus the unconditional variance which made it possible for volatility to become forecastable based on past shocks. The ARCH processes were also shown to be leptokurtic, with the fat tails even when the shocks are Gaussian. And to check for the reliability of the ARCH model and whether it has unbiased results, Engel introduced the ARCH-LM test which still remains a standard diagnostic tool in econometrics.

In spite of the new pathways the ARCH approach introduce for empirical finance, practical estimation still becomes cumbersome because capturing long memory requires many lags(ARCH(q) with large \( q \)), making estimation of these long memory models inefficient.

\subsection{GARCH and IGARCH}
To allow the conditional variance to depend on both the lagged squared residuals and lagged conditional variance terms, \cite{bollerslev1986garch} introduced the Generalized ARCH (GARCH) model, which addressed the practical limitations of high order ARCH. Thus, the GARCH(1,1) quickly became the primary model of volatility modelling due to its highly parsimonious and empirical effectiveness.

The GARCH model, therefore, brought several advances. First, the GARCH(1,1) model often captures long-range dependencies with only two parameters, \(\alpha\) and \(\beta\), addressing parsimony. To address volatility persistence, the model use the sum of \(\alpha + \beta\) to quantify the speed of mean reversion, and the model shows that unconditional distributions remain fat-tailed.

Subsequently, \cite{englebollerslev1986persistence} further extended the theory by introducing the Integrated GARCH (IGARCH) model, in which the persistence parameter \(\alpha + \beta = 1\). In this specification, instead of volatility shocks decaying over time, they are allowed to accumulate indefinitely, implying that conditional variance behaves like 
a unit-root process. Ideally, IGARCH captures environments in which markets 
do not revert quickly to long-run volatility levels, and where new information, including small shocks, produces long-lasting effects on uncertainty.

The development of the IGARCH framework was particularly important for modelling exchange rates in the 1980s, where persistent volatility and slow mean reversion were empirically pervasive. 
In subsequent years, researchers found that cryptocurrencies exhibit remarkably similar characteristics: 
high persistence, weak mean reversion, and volatility that remains elevated long after major shocks. IGARCH, thus, serves as a bridge between classical GARCH methodologies and the empirical reality of digital asset markets, in which volatility dynamics often resemble near-unit-root behaviour.

\subsection{Continuous-time links}
To bridge the discrete-time GARCH with the continuous-time stochastic volatility models, \cite{nelson1990archdiffusion} introduced the exponential GARCH(EGARCH), showing that as the sampling intervals shrink, many GARCH processes converge to diffusion processes. This result provided the theoretical justification for using GARCH in contexts traditionally dominated by continuous-time models, such as option pricing. Nelson proved that GARCH(1,1) converges to a process with an inverted-gamma stationary distribution, while exponential-GARCH processes yield log-normal variance diffusions, thus providing more insights into the theoretical legitimacy for applying GARCH to high-frequency data and enriched the versatility of volatility models.

\subsection{Stylised facts}

\cite{bollerslev1994archsurvey} defined a set of stylised facts about financial return series after surveying more than a decade of empirical research, showing that these returns remain stable across assets, markets and time periods. These findings still serve as a guide for the specification, estimation and evaluation of volatility models in the GARCH tradition and beyond. 

Volatility clustering is an essential fact associated with financial asset returns as large price movements tend to be followed by large movements, either negative or positive, while periods of stability and peace tend to persist. This temporal dependence in the conditional variance tends to be one of the strongest characteristics of financial returns and forms the core empirical motivation for GARCH models as it suggests that volatility is persistent and forecastable, contradicting the notion of constant variance assumed in early asset pricing models prior to the introduction of the ARCH framework.
 
The authors identified that leverage effects are integral in financial markets since in these markets, volatility responds asymmetrically to news and bad news, in particular, with negative returns typically increases volatility more than good news. The authors attribute this phenomenon to financial leverage, highlighting that a fall in equity value increases leverage and raises risk, or to behavioural factors such as asymmetric reactions to losses. Thus, leverage effects are a defining feature in equities, they are notably weaker or even absent in Bitcoin and many cryptocurrencies, highlighting the fundamental differences in market microstructure.

Bollerslev and the co-authors found that empirical return distributions exhibit  kurtosis far exceeding that of the normal distribution meaning that the results are heavily fat-tailed and the financial returns display more frequent extreme movements than the Gaussian models predict. GARCH processes with conditional normal errors can already generate heavy-tailed unconditional distributions, a feature that is further amplified when using fat-tailed innovation distributions like the Student-\(t\).

The financial markets are found to rely heavily on information flow, and they experience predictable spikes in volatility around scheduled events such as monetary policy announcements, labour market releases, or earnings reports. Whereas these predictable bursts do not imply that the effect of information is neither uniform nor random, they have motivated the introduction of models that explicitly incorporate exogenous information or calendar effects, and it is particularly relevant to cryptocurrencies, where regulatory announcements, exchange hacks, and protocol updates often generate disproportionate volatility.

\cite{bollerslev1994archsurvey} found the existence of volatility co-movements across assets, noting that the volatilities volatilities across different assets, sectors, or markets tend to move together, reflecting the presence of common shocks or shared macroeconomic risks necessitates the need for multivariate GARCH models such as CCC, DCC, and BEKK, and explains why systemic or contagion effects can emerge during periods of market stress. Co-movement is also observed across cryptocurrencies, which often respond simultaneously to global risk sentiment or major news events.


These stylised facts jointly form the conceptual backbone for modern volatility modelling, justifying the need for dynamic conditional variance models like GARCH, as well as more flexible nonlinear approaches such as regime-switching, smooth-transition models, realised volatility frameworks, and machine learning methods. The persistence, asymmetry, tail behaviour, and cross-asset dependence described by these facts are very persistent in cryptocurrency markets, making them crucial considerations in any attempt to model Bitcoin volatility.


\section{First generation cryptocurrency volatility studies (2016–2019)}
The initial research on cryptocurrencies applied the classical GARCH-family models with the aim of assessing whether Bitcoin behaved like traditional assets or not, identifying important empirical characteristics like extreme volatility, persistence patterns, and possible regime shifts which have served as defining features for the subsequent research in this area.

\subsection{Symmetry and volatility persistence}
To position the asset within the broader volatility landscape, \cite{dyhrberg2016bitcoin} carried out one of the earliest thorough GARCH applications to Bitcoin, comparing its behaviour to gold and the U.S. dollar. Using both GARCH(1,1) and EGARCH specifications, the study demonstrated several important features that subsequently shaped much of the early cryptocurrency volatility literature.

First, the author discovered high volatility persistence similar to commodity markets, implying that  shocks to Bitcoin variance disappear slowly and can influence volatility over extended periods. This is consistent with the later findings as is \cite{chu2017garchcrypto} and \cite{katsiampa2017volatility}, which frequently identified IGARCH-type behaviour for Bitcoin and other cryptocurrencies.

Second, the existence of symmetric volatility responses, indicating that positive and negative return shocks have similar effects on the volatility of Bitcoin, which sharply contrasts with equity markets, where negative shocks typically trigger larger volatility increases. In their subsequent study, such as \cite{bencheikh2020stgarch} confirmed this phenomenon and even documented \textit{inverted asymmetry} for certain cryptocurrencies, where volatility sometimes increases more after positive news, a feature contradicting the traditional leverage effect in financial markets.

Finally, since Bitcoin exhibited weak correlations with other equities, \cite{dyhrberg2016bitcoin} suggested that it can be used potentially as a diversifying asset and hedge against equities, thus assuming roles similar to the precious metals like gold. This sets the stage for exploring how Bitcoin as an asset behaves in relation to other assets as well as macroeconomic and behavioural variables.

Thus, from the foregoing, Dyhrberg’s results' implications are three-fold. First, the absence of leverage effects hints at Bitcoin’s distinctive behavioural and informational environment, characterised by retail-driven sentiment, speculative trading, and limited balance sheet mechanisms compared to equities, a fact consistent with later studies showing that sentiment and attention variables such as forum discussions as presented by
\cite{mai2018silentmajority}, news narratives by \cite{sapkota2022newsbitcoin} and \cite{coulter2022media}, and Google search intensity play a first-order role in driving volatility.

Second, the symmetry of volatility responses suggested early on that classical GARCH extensions designed for leverage effects (e.g., EGARCH, GJR-GARCH) might be less relevant for Bitcoin. Indeed, empirical comparisons in \cite{chu2017garchcrypto} and \cite{katsiampa2017volatility} often found that symmetric models (GARCH, IGARCH, CGARCH) performed as well or better than asymmetric ones.

Finally, Dyhrberg’s work established Bitcoin as an asset whose volatility behaviour is closer to commodities than equities, thereby reinforcing the need for models capable of capturing persistent, sentiment-sensitive, and regime-dependent dynamics which serves as the basis for later research using regime-switching GARCH models \citep{ardia2019regime}, smooth-transition 
models \citep{bencheikh2020stgarch}, and hybrid machine learning approaches 
\citep{zahidetal2022hybridgarch, han2025combinedmethod}, all of which attempt to account for nonlinearities and evolving market conditions not captured by early single-regime models.


\subsection{Comprehensive GARCH comparisons}
In carrying out one of the first large-scale, systematic comparisons of the GARCH-type models applied to digital assets, \cite{chu2017garchcrypto} estimated twelve volatility specifications, including GARCH, EGARCH, GJR-GARCH, FIGARCH, and IGARCH across seven major cryptocurrencies, with results revealing several distinctive features of cryptocurrency volatility dynamics.
In the first place, IGARCH frequently emerged as the best model for modelling the volatilities of cryptocurrencies, suggesting that volatility shocks in Bitcoin and other cryptocurrencies exhibit near-unit-root persistence, a finding consistent with the long-memory properties of Bitcoin which was later explored through component and regime-switching models.
Secondly, normal innovations were found to often outperform heavy-tailed alternatives, a somewhat surprising result given the strong empirical evidence of fat tails in traditional equity and commodity returns. This result may reflect the dominance of structural breaks, speculative bursts, and regime shifts in crypto markets rather than heavy-tailed shock distributions per se.
Finally, different cryptocurrencies were found to exhibit different volatility structures, with some assets displaying stronger asymmetry, longer memory, or higher sensitivity to shocks than others. This heterogeneity suggests that cryptocurrencies cannot be modelled as a uniform asset class given that they may reflect differences in liquidity, maturity, and investor composition, key factors which influence the performance of the returns of cryptocurrencies.

Thus, Chu and the co-authors positioned IGARCH as a natural benchmark for Bitcoin volatility modelling and highlighted the need for models capable of capturing long persistence and structural heterogeneity. In an indirect response to this study, subsequent research including component GARCH models 
by \cite{katsiampa2017volatility}, regime-switching GARCH frameworks by \cite{ardia2019regime}, and smooth-transition approaches as in \cite{bencheikh2020stgarch} bring more light to the empirical complexities that \cite{chu2017garchcrypto} identified.


\subsection{Component GARCH}

In comparing nine GARCH-type models including the standard GARCH, EGARCH, GJR-GARCH, TGARCH, and FIGARCH, \cite{katsiampa2017volatility} found that the Component GARCH (CGARCH) model offered the best overall fit among these GARCH-family models for cryptocurrencies. CGARCH decomposes volatility into a short-run transitory component and a long-run persistent component, thereby capturing both immediate reactions to market news and deeper structural trends. The finding that Bitcoin's volatility is driven by both short-lived shocks like speculative bursts, exchange incidents, regulatory announcements and slowly evolving forces for example, adoption rates, liquidity growth, macroeconomic uncertainty provides an important 
conceptual bridge between earlier research work.

In particular, the CGARCH superiority aligns closely with evidence from \cite{chu2017garchcrypto}, who found the existence extreme persistence and near-unit-root behaviour in cryptocurrencies, often favouring IGARCH specifications. Katsiampa's results refine this view by showing that what appears to be IGARCH-type persistence in single-regime models may actually 
reflect a combination of long-term variance dynamics and short-term speculative noise. This is a consistent interpretation as was the case in \cite{dyhrberg2016bitcoin}, who also observed persistent volatility and symmetric responses to shocks, suggesting that Bitcoin’s volatility is not purely driven by episodic events but reflects deeper market forces.

Thus, Katsiampa’s contribution serves as an intermediate step between the early single-regime GARCH evidence and the nonlinear, regime-switching models that followed in later studies, for example, \cite{ardia2019regime}. 

\subsection{Regime switching}

\cite{ardia2019regime} provided evidence that Bitcoin’s volatility dynamics are regime-dependent by identifying several low-volatility and high-volatility regimes, highlighting through Markov-switching GARCH models (MSGARCH) three key insights. First, regime switching explains the appearance of IGARCH-like persistence in single-regime models and secondly, high-volatility regimes coincide with market crises and speculative booms, and finally, Value-at-Risk forecasts improve under regime-switching models, findings which have motivated later research on nonlinear models and hybrid specifications.

\section{Asymmetry, smooth transitions, and realised volatility (2019–2022)}

Volatility in cryptocurrencies is often asymmetric, persistent, and subject to structural breaks, hence the literature has expanded to include models with nonlinear and smooth-transition dynamics, as well as realised volatility models that exploit high-frequency data.

\subsection{Asymmetric and smooth-transition dynamics}

\cite{bencheikh2020stgarch} introduced Smooth Transition GARCH (ST-GARCH) models to analyse gradual, rather than abrupt, shifts between low-volatility and high-volatility states in major cryptocurrencies and found that in the first place, Bitcoin, Ripple and Litecoin portray inverted asymmetry where volatility increases more after positive returns, contrasting strongly with the negative-return leverage effects observed in equities. They also found that volatility responses can vary across cryptocurrencies depending on maturity, liquidity and use base as depicted by Ethereum that exhibits negligible asymmetry. Finally, Bitcoin exhibited smoother transitions between volatility states than smaller cryptocurrencies which indicates a relatively more stable and liquid market structure.

This is a confirmation of the existence of symmetry in \cite{dyhrberg2016bitcoin} and the dominance of symmetric GARCH models in \cite{chu2017garchcrypto}, while extending the literature by showing that asymmetry in crypto markets is not only weak but, in some cases, reversed. The inverted asymmetry pattern supports interpretations of speculation, heavy reliance on momentum-driven trading behaviour, and reinforces the view that volatility 
in cryptocurrencies is shaped more by behavioural dynamics than by conventional 
leverage mechanisms.


\subsection{Realised volatility models}

Realised volatility models gained prominence in cryptocurrency research due to their ability to capture intraday information. \cite{corsi2009har} proposed the Heterogeneous Autoregressive (HAR-RV) model, which approximates long-memory volatility using daily, weekly, and monthly realised volatilities. In the subsequent studies applying HAR-RV to Bitcoin found that: HAR-RV captures long-memory effects without fractional integration and often, rivals or outperforms GARCH models in out-of-sample forecasts, concluding that high-frequency realised volatility reduces noise and improves forecast accuracy. Later,  \cite{dudek2023bitcoinrv} would show further proof that hybrid HAR-ML variants have more improvements on the standalone HAR-RV models.

\section{Behavioural and sentiment drivers (2018–2025)}
Since the cryptocurrency market is heavily influenced by online narratives, media cycles, social media activities and investor attention, behavioural finance plays an integral role in pricing and valuation of cryptocurrencies. This section, therefore, covers research on forums, news, Google trends, topic modelling, and sentiment–econometric hybrids.

\subsection{Social media and forums}

\cite{mai2018silentmajority} analysed BitcoinTalk forum posts and found that the sentiment from less active (“silent majority”) users strongly predicts next-day prices, implying that there are ``key" investors who directly influence the price of Bitcoin simply by what they say or post online. On the other hand, highly active users or the “vocal minority” influence trading volume, the actual raw numbers of Bitcoin units traded per day, and not the returns. The authors also show that forum sentiment provides stronger predictive power than Twitter sentiment, a sensible finding given that the `\textit{Bitcointalk}' that they study covers, almost exclusively, news on Bitcoin, unlike Twitter/X on which all sorts of items or news articles are shared.

Thus, Mia and the co-authors highlight the decentralised, retail-driven nature of cryptocurrency markets, where specific community opinions rather than large, influential individuals shape short-term price dynamics. This aligns with later evidence showing that crowd sentiment, search intensity, and collective attention often precede volatility spikes \citep{wang2023cryptoML}. This, therefore, provides a behavioural rationale for incorporating sentiment variables into volatility models by distinguishing between market participation intensity and predictive sentiments.


\subsection{News sentiment}

By extracting news articles from major English newspapers, \cite{sapkota2022newsbitcoin} found that mixed positive-negative sentiment has stronger effects on the direction of prices than purely positive or negative sentiment. This is because the mixed positive-negative is a clearer reflection of the real life trading situation. The author also found that sentiment, in general, significantly improves the HAR-RV forecasts, and weekly news is more useful than the daily news in the HAR-RV forecasting framework.

The emphasis on mixed-tone narratives suggests that ambiguity and conflicting news signals generate heightened uncertainty and speculative reactions among investors and the stronger predictive power at the weekly horizon indicates that sentiment diffusion in crypto markets may operate more slowly than in equity markets, possibly due to heterogeneous investor bases and global time-zone effects.

\subsection{Search-based attention}

Search-based measures capture the attention of retail investors and
\cite{hoangvo2024googleattention} identified that abnormal search volume (ASVI) is integral in predicting next-week crypto-returns, and also influences the increase in volatility and trading volume. ASVI is also shown to have stronger effects on larger cryptocurrencies in comparison  to the smaller units.

This reinforces the attention-based asset pricing literature and justifies the use of Google Trends as a timely behavioural proxy, especially in markets lacking traditional fundamental anchors. The finding that ASVI influences volatility is consistent with the “information-overload” hypothesis, where surges in retail attention precede market disorder 
and speculative trading. Later ML studies by \citep{wang2023cryptoML} incorporate Google Trends as an input feature, showing that attention-based variables enhance nonlinear models more effectively than standard GARCH specifications.


\subsection{Topic modelling and narrative effects}

\cite{coulter2022media} applied Latent Dirichlet Allocation (LDA) to global news coverage and identified three dominant narrative categories: crypto-crime, crypto-governance, and crypto-economy. Crypto-crime and governance narratives often preceded sharp daily declines in Bitcoin prices, especially during periods of regulatory uncertainty and exchange hacks, showing that volatility is not only driven by numerical sentiment indices but also by the thematic content and narrative framing of news, a finding consistent with 
\cite{sapkota2022newsbitcoin} who indicated that sentiment-driven volatility depends not only on tone but also on the type of narrative being amplified. 

\subsection{Hybrid sentiment -- econometric models}

Hybrid sentiment–econometric models integrate behavioural indicators with traditional time-series structures. For instance, \cite{kabo2025sentimenthybrid} showed that sentiment and GDP improve ARIMA and LSTM forecasts, highlighting the complementary value of macroeconomic and behavioural factors and \cite{leberegyed2025sentimentlstm} demonstrated that negative sentiment Granger-causes volatility and improves VaR forecasts when incorporated into GARCH–LSTM hybrids.

These studies emphasise that sentiment affects both central tendency and tail risk. They also show that hybrid models can extract nonlinear behavioural effects that GARCH models alone cannot capture, reinforcing the motivation for this study to compare traditional econometric and machine learning approaches in a sentiment-augmented framework.


\subsection{Random forests}

\cite{basher2022bitcoinRF} showed that Random Forest(RF) models outperform logistic regression in predicting Bitcoin price direction. This is, in part, because Random Forest models handle nonlinear interactions and high-dimensional data and are robust to multicollinearity and noise, hence provide interpretable feature importance measures.

RF has also been used to model realised volatility by \citep{luong2018harjl} and forecast volatility in emerging markets  by \citep{diane2024southafrica}, applications which indicate that tree-based models can successfully capture irregular patterns, structural breaks, and nonlinearities which features prevalent in Bitcoin markets. Thus, their flexibility makes them strong candidates for sentiment-augmented volatility forecasting.


\subsection{Deep learning}

Deep learning models such as LSTM, GRU, and CNN architectures outperform traditional econometric models in many studies. \cite{kimwon2018lstmhybrid} proposed LSTM–GARCH hybrids and subsequently, \cite{wang2023cryptoML} found that LSTM models outperform GARCH for weekly and  monthly volatility. Eventually, in an attempt to improve accuracy, \cite{han2025combinedmethod} used GARCH forecasts as LSTM inputs, and the result was a drastic improvement in the prediction accuracy of the LSTM model.

These improvements arise because deep learning architectures capture long-memory patterns, nonlinear relationships, and structural regime changes which are dynamics often persistent in Bitcoin volatility and amplified by sentiment shocks.


\subsection{Hybrid GARCH -- machine learning(ML) models}
In an attempt to improve the forecasting powers of the traditional econometric models, researchers have proposed hybrid models, combining the structural interpretability of these econometric models with the adaptability of the machine learning approaches. By assessing the GARCH residual learning, \citep{zahidetal2022hybridgarch} showed that ML models learn residual patterns not captured by GARCH. And to incorporate news sentiments and their effects on the model performance, \citep{leberegyed2025sentimentlstm} introduced the sentiment-augmented hybrids and reported that these behavioural variables improve performance. Finally, by combining GARCH-based features with graph neural networks, \citep{yu2025genshin} reported multi-scale volatility decomposition.

\subsection{Random forests for volatility}

Random Forests(RF) have been used not only for classification but also for volatility regression. \cite{amoanu2025arimacomparison} and \cite{demirer2022riskaversion} show that RF models can outperform econometric models for certain financial time series, especially when 
nonlinear predictors or behavioural indices are included. In realised volatility forecasting, RF improves over HAR-JL specifications by capturing nonlinear interactions \citep{luong2018harjl}, thus reinforcing the potential value of RF in modelling Bitcoin volatility, particularly when sentiment variables are incorporated.


\section{Research gap}

The above mentioned literature reviewed reveals several insights:
first, Bitcoin volatility is extraordinarily persistent, often approximating IGARCH behaviour and secondly, volatility dynamics are nonlinear, asymmetric, and regime-dependent. Third, sentiment which is captured in the reported literature through search activity, news tone, and forums, is a major driver of volatility. Fourth, machine learning and hybrid models consistently outperform standalone econometric approaches.

Yet, despite these clear indications, important gaps still remain in literature. Few studies directly compare GARCH-family models with Random Forests for volatility forecasting and sentiment is rarely integrated systematically across model classes for longer datasets spanning major post-2023 structural events are underutilised.

This dissertation addresses these gaps by conducting a sentiment-driven comparative analysis of GARCH and Random Forest models using Bitcoin data from 2017–2025.


\newpage
%===============================================================
\chapter{Methodology}
%===============================================================

\section{Introduction}

This chapter presents the econometric and machine learning methodology used to forecast Bitcoin volatility and to evaluate whether sentiment-enhanced nonlinear methods outperform standard GARCH benchmarks. The chapter specifies the data-generating process for returns, the conditional variance dynamics under GARCH-type models, the Random Forest forecasting set-up, the out-of-sample forecasting design, and the forecast evaluation framework and multiple-model comparison procedures.

\section{Notation and return construction}

Let $P_t$ denote the daily closing price of Bitcoin on day $t$. The continuously compounded (log) return is defined as
\begin{equation}
r_t = 100\left(\ln P_t - \ln P_{t-1}\right),
\end{equation}
where the factor $100$ expresses returns in percentage terms. Let $\mathcal{F}_{t-1}$ denote the information set available at time $t-1$. This specification follows standard practice in financial econometrics, where returns are modelled as conditionally heteroskedastic processes with time-varying variance as in \cite{bollerslev1986garch}.

A general conditional mean specification is
\begin{equation}
r_t = \mu + \sum_{i=1}^{p}\phi_i r_{t-i} + \varepsilon_t,
\end{equation}
where $\varepsilon_t$ is the innovation satisfying
\begin{equation}
\mathbb{E}(\varepsilon_t \mid \mathcal{F}_{t-1}) = 0, 
\qquad 
\mathbb{V}(\varepsilon_t \mid \mathcal{F}_{t-1}) = \sigma_t^2.
\end{equation}
The focus of the study is on modelling and forecasting $\sigma_t^2$.

\section{Realised volatility proxy}

Since true volatility is latent, evaluation is based on a realised volatility proxy constructed from returns. A standard daily proxy is the squared return,
\begin{equation}
RV_t = r_t^2,
\end{equation}
which is a noisy but widely used estimator of daily integrated variance when intraday data are unavailable. In robustness analysis, alternative proxies based on rolling windows may be considered, for example
\begin{equation}
RV^{(h)}_t = \frac{1}{h}\sum_{j=0}^{h-1} r_{t-j}^2,
\end{equation}
for horizon $h \in \{5,10,20\}$.
Although squared returns are a noisy proxy for latent volatility, they are widely used in daily-frequency studies where intraday data are unavailable \citep{andersen1998answer, patton2011volatility}.

\section{Econometric volatility models}

\subsection{Shock decomposition and innovation distribution}

Innovations are written as
\begin{equation}
\varepsilon_t = \sigma_t z_t,
\end{equation}
where $\{z_t\}$ is i.i.d. with $\mathbb{E}(z_t)=0$ and $\mathbb{V}(z_t)=1$. To accommodate fat tails commonly observed in cryptocurrency returns, estimation is conducted under alternative distributional assumptions for $z_t$, including the Gaussian and Student-$t$ distributions:
\begin{equation}
z_t \sim \mathcal{N}(0,1)
\quad \text{or} \quad
z_t \sim t_\nu \ \text{scaled to unit variance}.
\end{equation}
Allowing for Student-$t$ innovations accommodates the excess kurtosis commonly observed in cryptocurrency returns and is standard in volatility modelling \citep{hansen1994autoregressive}.

\subsection{GARCH(1,1)}

The baseline conditional variance specification is GARCH(1,1):
\begin{equation}
\sigma_t^2 = \omega + \alpha \varepsilon_{t-1}^2 + \beta \sigma_{t-1}^2,
\end{equation}
with parameter restrictions
\begin{equation}
\omega>0,\qquad \alpha\ge 0,\qquad \beta \ge 0.
\end{equation}
A sufficient condition for covariance stationarity is
\begin{equation}
\alpha+\beta < 1,
\end{equation}
in which case the unconditional variance exists and equals
\begin{equation}
\mathbb{E}(\sigma_t^2) = \frac{\omega}{1-\alpha-\beta}.
\end{equation}
The persistence of volatility shocks is measured by $\alpha+\beta$, and the corresponding half-life can be reported as
\begin{equation}
HL = \frac{\ln(0.5)}{\ln(\alpha+\beta)}.
\end{equation}
%The GARCH(1,1) model is widely used as a benchmark in volatility forecasting due to its parsimonious structure and strong empirical performance across asset classes \citep{bollerslev1986garch}.


\subsection{IGARCH(1,1)}

To allow near-unit-root volatility behaviour, the IGARCH(1,1) model imposes
\begin{equation}
\alpha+\beta = 1,
\end{equation}
so that
\begin{equation}
\sigma_t^2 = \omega + \alpha \varepsilon_{t-1}^2 + (1-\alpha)\sigma_{t-1}^2.
\end{equation}
Under IGARCH, shocks to volatility decay very slowly, which may reflect the high persistence documented in Bitcoin volatility. Near-unit-root volatility behaviour has been documented in cryptocurrency markets, motivating the use of IGARCH-type specifications to capture highly persistent volatility dynamics. %\citep{chu2017garchcrypto}.
\subsection{Sentiment-augmented variance dynamics}

To test whether investor attention contains incremental information for volatility dynamics, sentiment is incorporated via an exogenous regressor $S_{t-1}$ (Google Trends index or its transformation). A parsimonious sentiment-augmented specification is
\begin{equation}
\sigma_t^2 = \omega + \alpha \varepsilon_{t-1}^2 + \beta \sigma_{t-1}^2 + \gamma S_{t-1},
\end{equation}
where $\gamma$ captures the marginal contribution of attention shocks to next-day conditional variance. Nonlinearities can be partially addressed by including transformations such as $\Delta S_{t-1}$ or $\ln(1+S_{t-1})$.
%This approach is consistent with studies showing that investor attention and information demand contain predictive power for volatility and trading activity, particularly in cryptocurrency markets \citep{da2011attention, bucci2021stacking}.


\section{Estimation via maximum likelihood}

Let $\theta$ denote the vector of parameters (including $\omega,\alpha,\beta$ and distributional parameters such as $\nu$ for Student-$t$). Conditional on $\mathcal{F}_{t-1}$, the log-likelihood contribution is
\begin{equation}
\ell_t(\theta) = \log f\!\left(z_t(\theta)\right) - \frac{1}{2}\log(\sigma_t^2(\theta)),
\qquad
z_t(\theta)=\frac{\varepsilon_t}{\sigma_t(\theta)},
\end{equation}
where $f(\cdot)$ is the density of $z_t$. The sample log-likelihood is
\begin{equation}
\mathcal{L}(\theta)=\sum_{t=1}^{T}\ell_t(\theta),
\end{equation}
and $\widehat{\theta}$ is obtained by maximising $\mathcal{L}(\theta)$ subject to parameter restrictions.
Maximum likelihood estimation used since it is standard for GARCH-type models and yields consistent and asymptotically normal estimators under regularity conditions \citep{bollerslev1986garch}.

\section{Forecast definition}

A one-step-ahead conditional variance forecast from an estimated GARCH-type model is
\begin{equation}
\widehat{\sigma}_{t+1|t}^2
=
\widehat{\omega} + \widehat{\alpha}\widehat{\varepsilon}_{t}^2 + \widehat{\beta}\widehat{\sigma}_{t}^2
\quad
(\text{plus } \widehat{\gamma} S_t \text{ if sentiment is included}).
\end{equation}
These forecasts are generated recursively in a rolling-window design to mimic real-time updating.

\section{Random Forest forecasting set-up}

Let $\mathbf{X}_t$ denote the feature vector containing lagged returns, volumes, rolling volatility measures, and sentiment lags. The Random Forest produces a forecast
\begin{equation}
\widehat{RV}_{t+1|t} = f(\mathbf{X}_t),
\end{equation}
where $f(\cdot)$ is estimated as an ensemble of regression trees. To align targets across frameworks, machine learning outputs are mapped to a variance forecast and compared to realised volatility proxies. Hyperparameters are then selected using time-series cross-validation.

\section{Forecast evaluation: loss functions and model comparison}
Forecast evaluation follows the framework proposed by \cite{patton2011volatility}, which emphasises the use of loss functions that are robust to latent volatility and forecast dependence. In particular, the QLIKE loss function is employed due to its desirable theoretical properties when volatility is unobserved and must be proxied. 
\\Model comparisons are conducted using the model confidence set procedure of \cite{hansen2011model} and the superior predictive ability test of \cite{hansen2005test}, which jointly account for serial correlation, multiple testing, and data-snooping concerns.

\subsection{QLIKE loss}

Following Patton-style volatility evaluation, the primary loss function is QLIKE:
\begin{equation}
L_{t+1}(m) = \frac{RV_{t+1}}{\widehat{\sigma}^2_{t+1|t}(m)} 
- \ln\!\left(\frac{RV_{t+1}}{\widehat{\sigma}^2_{t+1|t}(m)}\right) - 1,
\end{equation}
where $\widehat{\sigma}^2_{t+1|t}(m)$ denotes the forecast from model $m$.

\subsection{Loss differentials and benchmark comparison}

Let model $0$ denote the benchmark (e.g., GARCH(1,1)). Define the loss differential for model $m$ relative to the benchmark as
\begin{equation}
d_{t+1}(m) = L_{t+1}(0) - L_{t+1}(m).
\end{equation}
A positive mean of $d_{t+1}(m)$ indicates that model $m$ improves upon the benchmark on average.

\subsection{Superior predictive ability (SPA) test}

The SPA test evaluates whether any competing model outperforms the benchmark while controlling for data-snooping. The null hypothesis is
\begin{equation}
H_0:\ \max_{m \in \mathcal{M}} \mathbb{E}\left[d_{t}(m)\right] \le 0,
\end{equation}
against the alternative that at least one model has strictly positive expected improvement. Inference is conducted using a block bootstrap to account for serial dependence in the loss differentials.

\subsection{Model confidence set (MCS)}

To compare multiple models simultaneously, the model confidence set procedure is used to identify a subset $\widehat{\mathcal{M}}_{1-\alpha}$ of models that are statistically indistinguishable from the best model under the chosen loss function. The MCS algorithm iteratively removes inferior models based on tests of equal predictive ability, using bootstrap-based critical values to accommodate dependence.

%===============================================================
%Chapter 3 with the math
% Comment out or Uncomment: Highlight, then Ctrl + Shift + 7
% \chapter{Methodology}
% %===============================================================

% \section{Introduction}

% This chapter outlines the methodological framework used to compare traditional 
% econometric models and machine learning techniques for forecasting Bitcoin 
% volatility. The focus is on clearly describing the models, data handling procedures, 
% sentiment integration, and the evaluation tools applied to assess forecasting 
% performance.

% \section{Research design}

% The study adopts a quantitative, comparative research design. Volatility forecasts 
% generated by GARCH-type models and Random Forests are evaluated against actual 
% realised volatility computed from daily Bitcoin returns. The analysis uses an 
% out-of-sample forecasting framework to compare the influence of sentiment-enhanced 
% features relative to traditional return-based predictors. The design allows for:
% \begin{itemize}
%     \item Model comparison under identical data and forecasting conditions.
%     \item Assessment of whether sentiment improves forecast accuracy.
%     \item Evaluation of whether nonlinear machine learning models outperform 
%     classical econometric models.
% \end{itemize}

% \section{Data description}

% The dataset consists of:
% \begin{itemize}
%     \item \textbf{Daily Bitcoin price series} (open, close, high, low, and trading volume), 
%     sourced from Yahoo Finance through the \texttt{yfinance} library.
%     \item \textbf{Google Trends sentiment proxy}, measuring search interest for the term 
%     ``Bitcoin,'' retrieved through the \texttt{pytrends} API.
% \end{itemize}

% Daily log-returns are computed from closing prices as:
% \begin{equation}
%     r_t = 100 \left( \ln P_t - \ln P_{t-1} \right),
% \end{equation}
% where $P_t$ is the Bitcoin closing price at time $t$.

% The sample covers \textbf{January 1, 2017 to September 25, 2025}, a period that includes 
% major price cycles, regulatory shifts, and structural changes such as the approval of spot 
% Bitcoin ETFs in 2024. Using this extended horizon will ensure that the models are 
% evaluated under diverse market conditions.

% \section{Data preprocessing}

% Several steps are applied to prepare the dataset for modelling:
% \begin{itemize}
%     \item \textbf{Return Computation}: Log returns are calculated from closing prices to 
%     stabilise variance and reduce scale dependence.
%     \item \textbf{Missing Data}: Both Bitcoin price and Google Trends data contains no missing values.
%     \item \textbf{Outlier Handling}: Extreme volatility spikes are retained because they 
%     represent genuine market events important for forecasting.
%     \item \textbf{Feature Engineering}: Predictor variables include lagged returns, lagged 
%     volatility, lagged sentiment indices, and past volumes. Multiple lags are used to 
%     capture short-memory and medium-memory effects highlighted in the literature.
%     \item \textbf{Stationarity Checks}: Tests (ADF, KPSS) are performed to confirm that the 
%     return series behaves consistently with standard assumptions for volatility modelling.
% \end{itemize}

% Returns follow the standard decomposition:
% \begin{equation}
%     r_t = \mu_t + \varepsilon_t, \qquad
%     \varepsilon_t = \sigma_t z_t,
% \end{equation}
% where $\sigma_t^2$ is the conditional variance and $z_t$ are i.i.d.\ innovations with
% $\mathbb{E}[z_t]=0$ and $\mathbb{Var}(z_t)=1$.

% Squared returns serve as the realised volatility proxy:
% \begin{equation}
%     v_t = r_t^{2}.
% \end{equation}

% \section{Econometric models}

% The study uses a set of GARCH-family models commonly applied in cryptocurrency 
% volatility research:
% \begin{itemize}
%     \item \textbf{GARCH(1,1)}: The classical benchmark model capturing volatility clustering.
%     \item \textbf{IGARCH(1,1)}: Incorporates near-unit-root behaviour, a feature frequently 
%     observed in Bitcoin volatility \citep{chu2017garchcrypto}.
% \end{itemize}

% The conditional mean is specified as:
% \begin{equation}
%     r_t = \mu + \varepsilon_t.
% \end{equation}

% The GARCH(1,1) conditional variance is:
% \begin{equation}
%     \sigma_t^2 = \omega + \alpha\,\varepsilon_{t-1}^2 + \beta\,\sigma_{t-1}^2,
% \end{equation}
% with constraints $\omega>0$, $\alpha\ge 0$, $\beta\ge 0$, and covariance stationarity
% when $\alpha + \beta < 1$.

% The IGARCH(1,1) specification restricts the persistence to unity:
% \begin{equation}
%     \sigma_t^2 = \omega + \alpha\,\varepsilon_{t-1}^2 + (1-\alpha)\sigma_{t-1}^2,
% \end{equation}
% implying:
% \begin{equation}
%     \alpha + \beta = 1.
% \end{equation}

% To incorporate sentiment, the variance equation is extended as:
% \begin{equation}
%     \sigma_t^2
%     = \omega
%     + \alpha\,\varepsilon_{t-1}^2
%     + \beta\,\sigma_{t-1}^2
%     + \gamma\,G_{t-1},
% \end{equation}
% where $G_{t}$ denotes the Google Trends index.

% These models are estimated via maximum likelihood under standard innovation 
% assumptions.

% \section{Random forest model}

% Random Forest (RF) is implemented as the primary machine learning approach for 
% volatility forecasting. RF is selected because:
% \begin{itemize}
%     \item It handles nonlinear relationships and interaction effects absent in GARCH.
%     \item It is robust to multicollinearity and noisy predictors.
%     \item It can incorporate sentiment features naturally.
%     \item It provides feature importance measures that facilitate interpretability.
% \end{itemize}

% The RF model uses:
% \begin{itemize}
%     \item Lagged returns and volumes,
%     \item Sentiment indices and their lags,
%     \item Rolling-window features such as moving averages of volatility.
% \end{itemize}

% Let the feature vector at time $t$ be:
% \begin{equation}
%     \mathbf{x}_t =
%     \big(
%         r_t, r_{t-1}, \dots,
%         V_t, V_{t-1}, \dots,
%         G_t, G_{t-1}, \dots
%     \big)^\top.
% \end{equation}

% The forecasting target is one-step-ahead volatility:
% \begin{equation}
%     y_{t+1} = r_{t+1}^2.
% \end{equation}

% A Random Forest consisting of $B$ trees generates forecasts as:
% \begin{equation}
%     \hat{y}_{t+1 \mid t}
%     = f_{\text{RF}}(\mathbf{x}_t)
%     = \frac{1}{B}\sum_{b=1}^{B} T_b(\mathbf{x}_t),
% \end{equation}
% where $T_b(\cdot)$ denotes the $b$-th regression tree.

% Appropriate hyperparameters (number of trees, depth, minimum samples per split) are 
% tuned using \textbf{time-series cross-validation} to avoid look-ahead bias.

% \section{Sentiment integration}

% Sentiment enters the models through Google Trends search volume, capturing public 
% interest and retail attention. The behavioural literature shows that attention shocks 
% predict volatility, returns, and trading activity. 

% This feature is included in both the econometric models (as exogenous regressors) 
% and machine learning models (as predictors), enabling a direct comparison of how 
% different modelling frameworks incorporate behavioural information.

% Formally:
% \begin{equation}
%     \text{Econometric models:}\quad \sigma_t^2 = f(\varepsilon_{t-1}^2,\sigma_{t-1}^2,G_{t-1}),
% \end{equation}
% \begin{equation}
%     \text{Random Forest:}\quad \mathbf{x}_t \supset \{G_t, G_{t-1}, \dots\}.
% \end{equation}

% \section{Forecasting framework}

% The forecasting design follows a rolling-window approach:
% \begin{itemize}
%     \item The models are trained using an initial in-sample period.
%     \item Forecasts are generated for the next day.
%     \item The window then rolls forward by one day, repeating the estimation–forecast cycle.
% \end{itemize}

% Let $W$ denote the estimation window length. For each forecast origin $t$:
% \begin{enumerate}
%     \item Estimate the model using data $\{r_{t-W+1}, \dots, r_t\}$.
%     \item Generate the one-step-ahead forecast $\hat{\sigma}_{t+1\mid t}^2$.
%     \item Update the window to $t \leftarrow t+1$.
% \end{enumerate}

% Compactly:
% \begin{equation}
%     \hat{\sigma}_{t+1\mid t}^2 
%     = \mathcal{M}\left(r_{t-W+1:t},\, X_{t-W+1:t}\right),
% \end{equation}
% where $\mathcal{M}$ is either a GARCH-type model or the Random Forest.

% \section{Evaluation metrics}

% To assess and compare model performance, the following evaluation tools are used:

% \subsection*{1. Root Mean Squared Error (RMSE)}
% Measures the average magnitude of forecast errors. RMSE penalises large errors more 
% heavily, making it suitable for detecting models that fail during extreme volatility.
% \begin{equation}
%     \text{RMSE}
%     = \sqrt{\frac{1}{T}
%         \sum_{t=1}^{T}
%         \left(\hat{\sigma}_t^2 - v_t\right)^2 }.
% \end{equation}

% \subsection*{2. Mean Absolute Percentage Error (MAPE)}
% Provides a scale-free error measure that is easy to interpret across time and models.
% \begin{equation}
%     \text{MAPE}
%     = \frac{100}{T}
%       \sum_{t=1}^{T}
%       \left|\frac{\hat{\sigma}_t^2 - v_t}{v_t}\right|.
% \end{equation}

% \subsection*{3. Mean Absolute Error (MAE)}
% Useful for assessing general accuracy without disproportionately penalising extreme 
% observations.
% \begin{equation}
%     \text{MAE}
%     = \frac{1}{T}
%       \sum_{t=1}^{T}
%       \left|\hat{\sigma}_t^2 - v_t\right|.
% \end{equation}

% \subsection*{4. Diebold--Mariano (DM) Test}
% Used to evaluate whether the difference in forecasting accuracy between two models is 
% statistically significant. This is crucial for determining whether sentiment-augmented 
% models outperform baseline versions.

% Let $d_t$ be the loss differential between models A and B:
% \begin{equation}
%     d_t = L\left(\hat{\sigma}_{t,\text{A}}^2,v_t\right)
%         - L\left(\hat{\sigma}_{t,\text{B}}^2,v_t\right),
% \end{equation}

% The DM statistic is:
% \begin{equation}
%     DM = \frac{\bar{d}}{\sqrt{\widehat{\text{Var}}(\bar{d})}},
% \end{equation}
% where $\bar{d}$ is the sample mean and $\widehat{\text{Var}}(\bar{d})$ is the 
% Newey--West HAC estimator.

% \section{Reasons for the choosing of these evaluation tools}

% The selected metrics cover both general accuracy (RMSE, MAE, MAPE) and formal 
% statistical testing (DM). Cryptocurrency volatility exhibits heavy tails and structural 
% breaks, so:
% \begin{itemize}
%     \item RMSE is ideal for capturing large forecast misses,
%     \item MAPE/MAE handle non-Gaussian errors well,
%     \item DM ensures that improvements from sentiment features or machine learning are 
%     not due to randomness.
% \end{itemize}

% \section{Summary}

% The methodology integrates traditional econometric approaches, machine learning 
% techniques, and behavioural indicators to analyse Bitcoin volatility. By maintaining a 
% consistent forecasting framework and using robust evaluation metrics, the study aims to provide 
% a clear comparison of whether sentiment and nonlinear modelling techniques improve 
% volatility predictions relative to classical GARCH methods.

%===============================================================

%===============================================================
\chapter{Empirical Results}
%===============================================================

\section{Descriptive statistics}
\subsection{Core distributional statistics}
\input{tables/table_4_1_core_descriptives}
\subsection{Quantiles and tail behaviour}
\input{tables/table_4_2_quantiles_tail_behaviour}
\subsection{Dependence and persistence diagnostics}
\includegraphics{figures/fig_4_1_3_acf_returns}
\includegraphics{figures/fig_4_1_3_acf_squared_returns}
\input{tables/table_4_3_ljung_box_tests}
\subsection{Normality and distributional tests}
\input{tables/table_4_4_jarque_bera_tests}
\subsection{Stationarity diagnostics}
\subsubsection{ADF tests}
\input{tables/table_4_5_adf_stationarity_tests}
\subsubsection{KPSS tests}
\input{tables/table_4_5_kpss_stationarity_tests}
\subsection{Correlation analysis}
\subsubsection{Correlation matrix}
\input{tables/table_4_6_correlation_matrix}
\subsubsection{Rank correlations}
\input{tables/table_4_6_correlation_matrix_spearman}

\section{Model estimation results}
\subsection{Model specification and estimation setup}
\input{tables/table_4_7_model_specifications_setup}
\subsection{Baseline econometric volatility models}
\subsubsection{GARCH(1,1) estimation function}
\input{tables/table_4_8_garch11_baseline_parameters}
\subsubsection{IGARCH(1,1) estimation parameters}
\input{tables/table_4_8_baseline_econometric_params_garch_igarch}
\subsection{Sentiment-augmented econometric models}
\input{tables/table_4_9_sentiment_augmented_econometric_params}
\subsection{Econometric models diagnostics}
\subsubsection{Diagnostics for baseline GARCH}
\input{tables/table_4_10_garch_residual_diagnostics}
\subsubsection{Diagnostics for GARCH-X}
\input{tables/table_4_11_garchx_residual_diagnostics}
\subsection{ARCH-LM test on standardized residuals}
\input{tables/table_4_12_arch_lm_tests}
\subsection{Baseline diagnostics for GARCH and IGARCH}
\input{tables/table_4_10_all_baseline_residual_diagnostics_garch_igarch}
\subsection{Sentiment augmented diagnostics for GARCH and IGARCH}
\input{tables/table_4_11_all_sentiment_residual_diagnostics_garchx_igarchx}
\subsection{The Random Forests model}
\subsubsection{In-sample fit diagnostics}
\input{tables/table_4_12_rf_in_sample_fit}
\subsubsection{Feature importance}
\input{tables/table_4_13_rf_feature_importance}

\section{Forecasting performance}
%\input{tables/forecasts}
\subsection{Econometric models forecasts}
Discuss
\subsection{Random forests forecasts}
Discuss
\subsection{Forecast evaluation and model comparison}
\input{tables/table_4_15_forecast_loss_summary}

\input{tables/table_4_17_forecast_loss_ranking}
%\section{Sentiment-augmented models}
\section{Robustness checks}
\subsection{Alternative forecast evaluation metrics}
ADD TABLE AND DISCUSS
%\input{tables/robustness_451_rmse_mae}
\subsection{Alternative volatility proxies}
%\input{tables/robustness_452_abs_returns}
ADD TABLE AND DISCUSS
\subsection{Alternative training and forecasting scheme - rolling window}
%\input{tables/robustness_453_restricted_oos}
ADD TABLE AND DISCUSS
\subsection{Robustness of sentiment effects}
%\input{tables/robustness_454_sentiment_lag1}
ADD TABLE AND DISCUSS

%===============================================================
\chapter{Discussion}
%===============================================================
\section{Interpretation of findings}
\section{Comparison with existing literature}
\section{Practical implications}
\section{Limitations}


%===============================================================
\chapter{Conclusion}
%===============================================================

\section{Summary of key findings}
\section{Contributions to knowledge}
\section{Policy and practice implications}
\section{Future research directions}

%===============================================================

%---------------------------------------------------------------
\bibliographystyle{apacite}
\bibliography{references}
\appendix
\chapter{The declaration of the use of AI tools}

The thesis work was prepared with limited use of generative artificial intelligence(AI).
\begin{itemize}
    \item In the first place, \href{https://curre.helsinki.fi/chat/general}{CurreChat} was used to help define the structure or the outline of the document.
    \item \href{https://chatgpt.com/}{ChatGPT} was used to summarise some of the papers use as reference material, language refinement and structural editing during the drafting process.
\end{itemize}
All substantive parts of the thesis work including the formation of research questions, model selection and data collection were conducted by the author. No AI-generated data, models or empirical results have been used in this study.\\
As the author, I take full responsibility for the accuracy of the content and originality of this thesis work.


%\chapter{Additional Tables}

%\chapter{Supplementary Figures}

%\chapter{Robustness Checks}

\end{document}
